\documentclass[10pt,a4paper]{amsart}
\usepackage[margin=19mm]{geometry}
\usepackage{amsmath,amssymb,mathtools}
\usepackage[scr=rsfs,cal=euler]{mathalfa}
\usepackage{xfrac}
\usepackage[unicode,hidelinks,bookmarks=false]{hyperref}
\newcommand{\ie}{i.e.\ }
\newcommand{\cf}{cf.\ }
\newcommand{\resp}{respectively\ }
\newcommand{\Subsection}{\subsection}
\providecommand{\nobreakdash}{\nobreak\hbox{-}}
\setlength{\emergencystretch}{2em}
\multlinegap0pt
\usepackage[shortlabels]{enumitem}
\usepackage{enumerate}
\usepackage{amssymb}
\usepackage{float}
\usepackage{xcolor}
\usepackage{tikz}
\newtheorem{example}{Example}
\newtheorem{lemma}{Lemma}
\newtheorem{prop}{Proposition}
\newcommand\BBK{{\mathbb K}}
\newcommand\BBQ{{\mathbb Q}}
\newcommand\CA{{\mathscr A}}
\newcommand\CB{{\mathscr B}}
\newcommand\GL{\operatorname{GL}}
\title{On connected subgraph arrangements}
\author{Lorenzo Giordani, Tilman M\"oller, Paul M\"ucksch, Gerhard R\"ohrle}
\date{Source: arXiv:2502.18144v2 (CC BY 4.0)}
\begin{document}
\maketitle

\noindent\emph{Source and licence.} On connected subgraph arrangements. Version arXiv:2502.18144v2 (CC BY 4.0), distributed by arXiv under the Creative Commons Attribution 4.0 International licence, http://creativecommons.org/licenses/by/4.0/, as recorded in the arXiv licence metadata for this exact version and shown on the abstract page https://arxiv.org/abs/2502.18144v2. Full licence notice: \texttt{licenses/arxiv-2502.18144-CC-BY-4.0.txt}.\par\smallskip

\noindent\textit{Editorial scope.} Counted unit: the article's prop prop:gen (source0.tex lines 1218--1223). The article's complete original proof of it is delivered with it (source0.tex lines 1224--1236, one \texttt{proof} environment of 13 lines). No mathematics is written, repaired or re-derived: the statement and the proof are the article's own lines, in the article's own notation, and no retained line is substituted or re-flowed.  Retained uncounted as the source-local context the proof needs: lines 1161--1179; lines 1196--1206.  Nothing was omitted from any retained span, so the package carries no omission record.  No retained line contains a citation, so the wrapper carries no bibliography and no bibliography entry.  No retained line contains a figure environment, an image-inclusion command or drawing code, so no image is delivered and the package declares no asset dependency.  Editorial formatting only, in addition: an AMS \texttt{amsart} preamble that restates the article's own declarations and macros used by the retained lines, namely the environments \texttt{example}, \texttt{lemma}, \texttt{prop} on separate counters, together with the standard packages those lines need.  The retained environments print in this document as Example 1, Lemma 1, Proposition 1 in order of appearance, whereas the article prints the same items with the numbering of its own full document.  The complete native source of the article as distributed in the arXiv e-print for this exact version is bundled unchanged as \texttt{sources/173832/source0.tex}.
\begin{example}%\hfill
	\label{ex:Proj(non)unique}
	%\begin{enumerate}[(1)]	\item 
	(i). Let $\CA$ be a generic arrangement in $V = \BBK^\ell$ with $|\CA| = \ell+1$.
			Then $\CA = \{H_1,\ldots,H_\ell,H'\}$ where $\alpha_1,\ldots,\alpha_\ell \in V^*$
			with $H_i = \ker(\alpha_i)$ and $\alpha_1,\ldots,\alpha_\ell$ form a basis of $V^*$.
			If $\gamma = \sum_{i=1}^\ell a_i\alpha_i \in V^*$ with $\ker(\gamma) = H'$, then by the genericity of $\CA$
			we have $a_i\neq 0$ for all $1\leq i \leq \ell$. Hence, after rescaling the $\alpha_i$,
			we see that $\CA$ is linearly isomorphic to $\{\ker(x_1),\ldots,\ker(x_\ell),\ker(\sum_{i=1}^\ell x_i)\}$,
			where $x_1,\ldots,x_\ell$ is some basis of $V^*$.
			Thus, $\CA$ is projectively unique.
			
		%\item 
		(ii). Let $V = \BBQ^2$ and $x,y$ be a basis of $V^*$.
			Let $\CA = \{\ker(x),\ker(y),\ker(x+y),\ker(x+2y)\}$ and $\CB = \{\ker(x),\ker(y),\ker(x+y),\ker(x+3y)\}$ be arrangements in $V$.
			Then clearly $\CA \cong_L \CB$ but $\CA$ is not linearly isomorphic to $\CB$.
			Thus, $\CA$ (resp.\ $\CB$) is not projectively unique.
	%\end{enumerate}
\end{example}

\begin{lemma}
	\label{lem:gen}
	Let $\CA$ and $\CB$ be two arrangements in $V$.
	Suppose $\varnothing \ne S \subseteq \CA$, $\varnothing \ne T \subseteq \CB$
	such that $\langle S \rangle_\CA = \CA$ and $\langle T \rangle_\CB = \CB$, i.e., $S$ generates $\CA$
	and $T$ generates $\CB$.
	If $\CA$ and $\CB$ are $L$-equivalent via a poset isomorphism $\psi:L(\CA)\to L(\CB)$ and $\varphi \in \GL(V)$
	such that $\psi(S) = \varphi(S) = T$ and $\psi(H) = \varphi(H)$ for all $H \in S$, 
	then $\varphi$ extends to a linear isomorphism between the whole arrangements, i.e., 
	$\varphi(\CA) = \CB$.
\end{lemma}

\begin{prop}
	\label{prop:gen}
	Let $\CA$ be an essential and irreducible arrangement in $V \cong \BBK^\ell$.
	Suppose there is a subset $S$ of $\CA$ such that $\langle S \rangle_\CA = \CA$ and $|S| = \ell+1$.
	Then $\CA$ is projectively unique.
\end{prop}

\begin{proof}
	Since $\CA$ is irreducible, after an appropriate coordinate change we may assume without loss of generality that $S = \{H_1=\ker(x_1),\ldots,H_\ell = \ker(x_\ell),H_{\ell+1}=\ker(\sum_{i=1}^\ell x_i)\}$
	where $x_1,\ldots,x_\ell$ is a basis for $V^*$ (cf.\ Example \ref{ex:Proj(non)unique}(i)).
	Let $\CB$ be another arrangement in $V$ which is $L$-equivalent to $\CA$ via $\psi:L(\CA)\to L(\CB)$ and set $T := \psi(S)$.
	Since $\psi$ is a poset isomorphism, we have $\langle T \rangle_\CB = \CB$.
	Moreover, $T$ spans a generic subarrangement of $\CB$ with $\ell+1$ hyperplanes, 
	say $T = \{\psi(H_1)= \ker(\beta_1),\ldots,\psi(H_\ell) = \ker(\beta_\ell),\psi(H_{\ell+1}) = \ker(\gamma)\}$.
	After rescaling $\beta_1,\ldots,\beta_\ell$ by suitable scalars, we can assume $\gamma=\sum_{i=1}^\ell\beta_i$
	and note that $\beta_1,\ldots,\beta_\ell$ yield a basis for $V^*$. Thus, if we take $\varphi \in \GL(V)$ 
	as the unique linear map which is dual to the transformation of $V^*$ which maps $x_1,\ldots,x_\ell$ to $\beta_1,\ldots,\beta_\ell$,
	we have $\psi(S) = \varphi(S) = T$ and $\varphi(H) = \psi(H)$ for all $H\in S$. 
	By Lemma \ref{lem:gen}, $\CB$ is linearly isomorphic to $\CA$ which concludes the proof.
\end{proof}

\end{document}
